% =====================================================================
\section{Introduction}\label{sec:intro}
\IEEEPARstart{L}{arge} language models (LLMs) are now routinely embedded in software engineering workflows~\cite{hou2024llm4se}. Frameworks such as MetaGPT~\cite{metagpt}, ChatDev~\cite{chatdev} and AutoGen~\cite{autogen} organize several LLM agents into roles (architect, programmer, tester) that hand artifacts to one another, and stateful orchestrators such as LangGraph~\cite{langgraph} wire such agents into graphs. Surveys identify LLM-based multi-agent systems as a central direction for automated software engineering~\cite{he2025mas,wang2024agents}.

Most quality assurance for these systems evaluates individual answers: is the generated code correct~\cite{evalplus,swebench}, does it parse~\cite{synchromesh,outlines}, does it pass a guardrail~\cite{nemoguardrails}? Multi-agent systems also fail \emph{between} answers. An empirical study of multi-agent LLM systems attributes many failures to inter-agent misalignment and to missing or incorrect verification of task completion~\cite{cemri2025mast}. Three situations illustrate failures that no per-answer check can detect:
\begin{itemize}
\item \emph{Stale answer.} A coding agent is asked to implement revision~1 of a requirement. While it works, the requirement changes to revision~2. The old answer arrives later, parses, passes its tests and is accepted. It answers a question that no longer exists.
\item \emph{Local validity, global conflict.} One design agent declares document identifiers as strings and another, for a neighboring service, as integers. Each output passes its own validator, yet the workspace is inconsistent.
\item \emph{Silent reconfiguration.} A security reviewer is added after the workflow looks complete. Unless completion is recomputed, the earlier ``done'' verdict is reported for a team that has not reviewed anything.
\end{itemize}

\textbf{Coordination is not a new problem.} These are coordination failures: they concern which inputs an answer was produced from, who was allowed to write which artifact, and when a collection of artifacts may be released. Multi-agent systems (MAS) research has studied coordination for more than three decades~\cite{malone1994coordination,jennings1993commitments,omicini2001coordination}. Coordination models such as Linda and programmable tuple centres mediate interaction through a shared medium whose reactions are specified outside the agents~\cite{gelernter1985linda,gelernter1992coordination,omicini2001tuple}. The Agents \& Artifacts (A\&A) meta-model makes the shared environment a first-class set of artifacts that agents use through operations, with the artifact enforcing its own guards~\cite{omicini2008artifacts,weyns2007environment,ricci2011environment}. Commitment-based and information-based protocols coordinate autonomous, possibly untrusted agents without inspecting their internal states~\cite{singh1998acl,yolum2002commitments,singh2011bspl,singh2020clouseau}, and electronic institutions use runtime governors that admit an agent's action only if it is permitted in the current institutional state~\cite{esteva2001institutions,arcos2005institutions}. Our controller belongs to this lineage: it is a coordination artifact, in the sense of A\&A, that plays the role of an institutional governor for a workflow of generated artifacts. More recently, database ideas have been carried over to LLM agents. A position paper argues that multi-agent systems should prioritize concurrency control~\cite{yang2026concurrency}. Concurrent proposals apply sagas, optimistic concurrency control or dependency-scoped revalidation to agent actions, plans and memories~\cite{chang2025sagallm,lyu2026coagent,chen2026planfence,li2026memtx}. Section~\ref{sec:related} compares them in detail. Three differences matter. In our design, the staleness judgment stays in the trusted coordinator and covers everything that shapes the question, including the instruction and the validator, not only the data an agent read. Invalidation is content-based with early cutoff. And retries are bounded per dependency episode, with proved terminal outcomes.

\textbf{What changes with generative agents.} LLM agents differ from the agents these models were designed for in three ways that matter for coordination. First, their contributions are \emph{generated artifacts} whose validity is relative to what the agent was shown: a well-formed reply that complies with every interaction protocol can still be wrong because its inputs have changed. The question to enforce is therefore not only ``is this action permitted now?'' but ``was this value produced from the inputs that are current now?''. Second, generation is costly and nondeterministic. After a change the coordinator must decide exactly which artifacts to regenerate, which is the problem incremental build systems solve with content-based dependency keys and early cutoff~\cite{mokhov,erdweg2015pluto}. Third, an LLM agent has no inspectable mental state and may return anything, late, twice or never. Correctness of the coordination therefore cannot depend on agent reasoning, only on checks over what the agents return.

\textbf{Approach.} We treat agents as \emph{untrusted proposers}. A deterministic controller owns structure, dependencies, acceptance and publication; agents only propose primitive values (text, code, JSON) for slots that the controller has created. For every slot, the controller records an exact \emph{dependency key} (in build-system terms, a content-based rebuild key) over the slot's read set, the shared context selected for it, and the contract (instruction, schema, validator, owner) under which it is generated. A proposal commits only if, at commit time, it answers the request that is still current, comes from the authorized owner, passes its validator, and leaves the artifact structure well formed; this is optimistic concurrency control~\cite{occ}. A change invalidates exactly the slots whose read sets changed, a per-episode budget bounds generation attempts, and a publication barrier releases only snapshots in which nothing is stale, nothing is pending and all joint constraints hold.

\textbf{Scope and division of labor.} The approach targets \emph{artifact-centric, feed-forward} workflows whose structure can be declared: model-driven pipelines, schema-driven configuration and interface synthesis, compliance documentation, and CI-gated generation of code and tests for a declared module structure. It does not target open-ended repository refactoring in which agents must invent new modules and cross-file structure; Section~\ref{sec:limitations} discusses how structural proposals can be admitted through an untrusted ingress step. We distinguish two kinds of guarantee. \emph{Coordination guarantees} (no stale or unauthorized commit, exact invalidation, bounded retries, honest completion) hold for any agent behavior and for any validator, however weak. \emph{Behavioral correctness} of the released artifacts is exactly as strong as the declared validators, and no stronger (Corollary~\ref{cor:domain}). We also distinguish who writes what (Section~\ref{sec:roles}): the generic controller and its TLA$^+$ model are written and verified once by framework builders; a workflow author declares stages, owners, read sets and validators in ordinary code, and never writes TLA$^+$.

\textbf{Contributions.} This paper makes the following contributions.
\begin{enumerate}
\item \emph{A coordination semantics for generative proposals} (Sections~\ref{sec:overview} and~\ref{sec:semantics}), positioned as artifact-mediated coordination. It combines a deterministic structural kernel, exact dependency keys and a guarded transition system, and it distinguishes a value that was \emph{accepted}, one that is currently \emph{usable}, and a snapshot that is \emph{publishable}.
\item \emph{Guarantees with explicit assumptions and counterexamples} (Section~\ref{sec:guarantees}). Four results carry the argument. \emph{Safety}: no stale, unauthorized or invalid value becomes usable (Theorem~\ref{thm:safety}). \emph{Single-episode repair}: with an upstream usability gate, a change starts at most one episode per affected job, whatever the scheduling (Theorem~\ref{thm:single}). \emph{Terminal outcome}: for acyclic workflows, every episode ends committed or failed, corrected in this revision for structural checks on suspended consumers (Theorem~\ref{thm:termination}). \emph{Refinement}: asynchronous execution refines atomic execution (Theorem~\ref{thm:refinement}). Supporting results cover structural integrity, exact invalidation, a generation bound, safe reconfiguration and bounded feedback epochs; several of them hold by construction, and we say so.
\item \emph{An executable instantiation} (Section~\ref{sec:instantiation}): a TLA$^+$ model with a TLC-checked refinement mapping and scoped TLAPS lemmas; a reference controller with an independent trace checker; and a generic kernel for arbitrary job graphs, differentially tested against the reference controller.
\item \emph{An evaluation} (Section~\ref{sec:evaluation}) with six research questions. It combines mutation-based guard-necessity checks, invalidation cost on independent and coupled workspaces, scheduling experiments, conformance scenarios and bounded model checking with two studies with real LLMs: \LlmModelsN{} open models coordinating a design--code--test workflow for document services, and \BenchModelsN{} models on the \BenchTasks{} tasks of HumanEval+ and MBPP+~\cite{evalplus}.
\end{enumerate}
All artifacts, every model call, raw logs and the scripts that generate every number in this paper are available (Section~\ref{sec:data}).

\textbf{Organization.} Section~\ref{sec:background} introduces the running example and the background, with a table relating our terms to their standard counterparts. Section~\ref{sec:overview} gives an overview, Section~\ref{sec:semantics} the formal semantics, and Section~\ref{sec:guarantees} the guarantees. Section~\ref{sec:instantiation} describes the instantiation and Section~\ref{sec:evaluation} the evaluation. Section~\ref{sec:threats} discusses threats to validity, Section~\ref{sec:related} related work, and Section~\ref{sec:limitations} limitations and future work.
